\SourceLeaf{167}ja Deine Arbeit bald bekannt. Gewiss wirst Du dafür den Dank zwar nicht
des großen Publikums ernten (zu dem auch mancher gehört, der für einen
geschickten Mathematiker gehalten wird), denn ich überzeuge
mich immer mehr davon, dass die Zahl wahrer Geometer äußerst gering ist und
die meisten die Schwierigkeiten solcher Arbeiten weder beurteilen
noch überhaupt verstehen können --- aber gewiss den Dank aller
derer, deren Urteil allein Dir wirklich etwas wert sein kann. --- In
Braunschweig ist ein Emigrant namens \textsc{Chauvelot}, ein
durchaus fähiger Geometer, der vorgibt, die Theorie der Parallellinien
vollständig begründet zu haben, und seine Arbeit demnächst drucken
lassen wird. Ich verspreche mir aber nichts von ihm. In
\textsc{Hindenburgs} Archiv, 9\textsuperscript{tes} Stück, befindet
sich ebenfalls ein neuer Versuch über denselben Gegenstand von
einem gewissen \textsc{Hauff}, der unter aller Kritik ist. .\, .\, .\, .\, .\, .\, .

\medskip
\begin{center}
\rule{0.3\textwidth}{0.4pt}
\end{center}
\medskip

\Needspace{6\baselineskip}
\begin{center}
\textsc{Gauss} an \textsc{Bolyai}. Braunschweig, 25. November 1804.
\end{center}

.\, .\, .\, .\, .\, .\, .\ Ich habe Deinen Aufsatz mit großem Interesse und
großer Aufmerksamkeit gelesen und mich sehr an dem echten, gründlichen
Scharfsinn erfreut. Du willst aber nicht mein leeres Lob, das auch
in gewisser Hinsicht schon deshalb parteiisch scheinen könnte, weil Dein
Gedankengang meinem sehr ähnlich ist, mit dem ich
einst den Gordischen Knoten zu lösen versuchte und bis jetzt vergeblich
zu lösen versucht habe. Du willst nur mein aufrichtiges, unverhohlenes
Urteil. Und das ist, dass mir Dein Verfahren noch
\textls{nicht} 
genügt. Ich will versuchen, den Stein des Anstoßes, den ich
noch darin finde (und der wiederum zu derselben Gruppe von Klippen
gehört, an denen meine Versuche bisher scheiterten), so klar
wie möglich ans Licht zu ziehen. Ich hoffe zwar noch
immer, dass jene Klippen einst, und noch vor meinem
Ende, eine Durchfahrt erlauben werden. Indessen habe ich jetzt so
manche andere Beschäftigung vor mir, dass ich gegenwärtig
daran nicht denken kann. Und glaube mir, es wird mich herzlich
herzlich freuen, wenn Du mir zuvorkommst und es Dir
